% Generated draft: requires mathematical and source review.
\documentclass[11pt]{article}
\usepackage[margin=25mm]{geometry}
\usepackage{fontspec}
\setmainfont{Latin Modern Roman}
\usepackage{amsmath,amssymb,mathtools,unicode-math}
\setmathfont{Latin Modern Math}
\usepackage{enumitem,needspace,xurl,hyperref}
\setlength{\parindent}{0pt}
\setlength{\parskip}{5pt}
\setlength{\emergencystretch}{3em}
\newcommand{\N}{\mathbb{N}}
\newcommand{\Z}{\mathbb{Z}}
\newcommand{\Q}{\mathbb{Q}}
\newcommand{\R}{\mathbb{R}}
\newcommand{\C}{\mathbb{C}}
\newcommand{\F}{\mathbb{F}}
\newcommand{\A}{\mathbb{A}}
\newcommand{\PP}{\mathbb{P}}
\newcommand{\E}{\mathbb{E}}
\newcommand{\eps}{\varepsilon}
\newcommand{\dd}{\,\mathrm{d}}
\begin{document}
\section*{UnsolvedMath ID: 30004828}
\textbf{Source problem code:} catalogue entry
\subsection{Definitions and mathematical statement}
Work over \(\C\), and write \(\mathbb N_0=\{0,1,2,\ldots\}\). Let \(G\) be a finitely generated group and \(\C G\) its group algebra. A left Yetter--Drinfeld module over \(\C G\) is a complex vector space \(V\) with a left \(G\)-action and a decomposition
\[
V=\bigoplus_{g\in G}V_g,
\qquad h\cdot V_g=V_{hgh^{-1}}\quad(g,h\in G).
\]
Its support \(\{g\in G:V_g\ne0\}\) is consequently a union of conjugacy classes. The category \({}^{\C G}_{\C G}\mathcal{YD}\) has linear maps preserving the action and grading as morphisms, and braiding
\[
c_{V,W}(v\otimes w)=g\cdot w\otimes v\qquad(v\in V_g).
\]
A Hopf algebra in this category has structure maps that are morphisms in the category, with the bialgebra compatibility interpreted using this braiding.

Consider a Hopf algebra object with a compatible nonnegative integer grading
\[
E=\bigoplus_{n\in\mathbb N_0}E_n
\]
in \({}^{\C G}_{\C G}\mathcal{YD}\). The integer grading here is distinct from the \(G\)-grading of a Yetter--Drinfeld module. Let \(C_n(E)\) denote the coradical filtration of its underlying coalgebra, whose zeroth term is the sum of its simple subcoalgebras. The requirements that \(E\) be connected and coradically graded mean
\[
C_0(E)=E_0=\C1,
\qquad C_n(E)=\bigoplus_{j=0}^{n}E_j\quad(n\in\mathbb N_0).
\]
No finite dimensionality of the homogeneous components, and no assumption that \(E\) is generated by \(E_1\), is imposed in the supplied question.

Write \(E\#\C G\) for the Radford biproduct, or bosonization, which is an ordinary Hopf algebra. Its underlying vector space is \(E\otimes\C G\), with multiplication
\[
(x\# g)(y\# h)=x(g\cdot y)\# gh.
\]
For a unital complex algebra \(A\), use the convention
\[
\operatorname{GKdim}A
=\sup_U\limsup_{m\to\infty}
\frac{\log\dim_{\C}(U^m)}{\log m},
\]
where \(U\) ranges over finite dimensional complex subspaces of \(A\) containing \(1\), and \(U^m\) is the span of products of \(m\) elements of \(U\). This convention also applies when finite generation of \(A\) is not assumed.

\textbf{Classification question (supplied Question 1).} For each finitely generated group \(G\), classify, or at least characterize, all connected, coradically graded Hopf algebra objects \(E\) as above satisfying
\[
\operatorname{GKdim}(E\#\C G)<\infty.
\]
The source does not specify the equivalence relation for the requested classification; in particular, it does not expressly say whether isomorphisms must preserve the integer grading in addition to the Yetter--Drinfeld Hopf algebra structure.

\textbf{Supplied scope qualification.} The original discussion states, invoking a result of Gromov, that one may assume \(G\) is nilpotent-by-finite, meaning that \(G\) has a nilpotent normal subgroup of finite index. The displayed question nevertheless quantifies over all finitely generated groups. The supplied literature assessment instead refers to finitely generated nilpotent groups. These scopes are recorded without silently replacing the displayed quantifier by either restriction.

\textbf{Related questions and context from the supplied excerpt.} For \(V\in{}^{\C G}_{\C G}\mathcal{YD}\), the Nichols algebra \(\mathcal B(V)\) is the quotient of the braided tensor Hopf algebra \(T(V)\), with \(V\) primitive, by its largest homogeneous braided Hopf ideal contained in \(\bigoplus_{n\ge2}T^n(V)\). The report states that, for \(V=E_1\) in Question 1, the subalgebra of \(E\) generated by \(V\) is isomorphic to \(\mathcal B(V)\), and describes \(E\) as a post-Nichols algebra of \(V\). Its related Question 2 asks to classify all such modules \(V\) satisfying
\[
\operatorname{GKdim}\bigl(\mathcal B(V)\#\C G\bigr)<\infty.
\]
This is a related target, not a replacement for Question 1. The excerpt also calls a graded connected braided Hopf algebra \(R\) generated by \(V\simeq R_1\) a pre-Nichols algebra of \(V\), and identifies \(\mathcal B(V)\) as the algebra that is simultaneously pre-Nichols and post-Nichols. It ends before giving the announced further reduction of Question 1 after Question 2; that missing continuation is not reconstructed here.
\subsection{Short English statement}
For every finitely generated group \(G\), describe the connected braided Hopf algebras over \(\C G\) whose grading is their coradical grading and whose bosonization has finite Gelfand--Kirillov dimension. The source allows a classification or a characterization. Its discussion of nilpotent-by-finite groups and its dated assessment concerning nilpotent groups have different scopes; neither is silently substituted for the stated domain. This is a statement-only draft requiring human mathematical review.
\subsection{Sources}
\begin{itemize}
\item Local UnsolvedMath registry, permanent ID \(30004828\), title \emph{Finite-Growth Hopf Algebras in Yetter--Drinfeld Categories}, exact source problem code \texttt{catalogue entry}, in the Oberwolfach Reports Open Problems collection.
\item Supplied report citation: \emph{Quantum Groups -- Algebra, Analysis and Category Theory} (2022), Oberwolfach Reports, pp.~2397--2458, DOI \texttt{10.4171/owr/2021/44}. Supplied source link and literature-source link labelled ``Oberwolfach Report 44'': \url{https://doi.org/10.4171/owr/2021/44}. The registry gives proposed year 2021; the supplied citation gives 2022. These dates are retained without resolving their relationship.
\item The supplied original excerpt identifies the contribution as \emph{On pointed Hopf algebras with finite GK-dimension}, by Nicol\'as Andruskiewitsch. It attributes the classification method to H.-J. Schneider and the speaker, and the nilpotent-by-finite reduction to a result of Gromov. The registry supplies no separate proposer. Its reference to a survey labelled [1] has no accompanying bibliographic entry in the supplied record.
\item Supplied literature source labelled ``Yang on finite-GK pre-Nichols algebras'': \url{https://arxiv.org/abs/2504.12206}. No fuller bibliographic identification is supplied.
\item Supplied literature assessment checked on \textbf{2026-08-22}: the registry labels the problem \texttt{partially\_solved} and states that the finite-GK classification program is complete in important diagonal and selected nondiagonal cases, while the requested classification over finitely generated nilpotent groups is not complete. This is supplied dated curation derived from the OWR record, not an independently verified literature assessment. The record expressly requests independent verification of time-sensitive status and cited claims.
\item The registry labels the statement \texttt{corrected\_verified} and reports that the linked report was inspected and its numbered Question 1 transcribed as a self-contained statement. That verification claim is supplied provenance; no source visit or independent verification is asserted in this draft.
\end{itemize}
\end{document}
